\chapter{Introduction}

\section*{Equation classes and ellipticity}

The basic linear and quasilinear second-order equations are, respectively,

\dcref{ch1:1.1}\group{chapter-1}\source{gilbarg-trudinger:page-0014}
\begin{equation}\label{gt.eq1.1}
\tag{1.1}
Lu \equiv a^{ij}(x)D_{ij}u+b^i(x)D_i u+c(x)u=f(x),\qquad i,j=1,2,\dots,n,
\end{equation}

and

\dcref{ch1:1.2}\group{chapter-1}\source{gilbarg-trudinger:page-0014}
\begin{equation}\label{gt.eq1.2}
\tag{1.2}
Qu \equiv a^{ij}(x,u,Du)D_{ij}u+b(x,u,Du)=0.
\end{equation}

Here $Du=(D_1u,\dots,D_nu)$, $D_i u=\partial u/\partial x_i$, and
$D_{ij}u=\partial^2u/(\partial x_i\,\partial x_j)$; repeated indices are
summed. Either equation is elliptic when $[a^{ij}]$ is positive definite
throughout its domain of arguments. It is uniformly elliptic when the ratio
$\gamma$ of its largest to smallest eigenvalue is bounded.

The model linear equations are Laplace's equation

\dcref{ch1:1.3}\group{chapter-1}\source{gilbarg-trudinger:page-0014}
\begin{equation}\label{gt.eq1.3}
\Delta u = \sum D_{ii}u = 0
\end{equation}

and Poisson's equation $\Delta u=f$. The minimal-surface equation

\dcref{ch1:1.4}\group{chapter-1}\source{gilbarg-trudinger:page-0014}
\begin{equation}\label{gt.eq1.4}
\sum D_i\!\left(D_i u/(1+|Du|^2)^{1/2}\right)=0
\end{equation}

is quasilinear and non-uniformly elliptic, with
$\gamma=1+|Du|^2$.

\section*{Linear theory}

For Laplace's equation, the Perron construction of harmonic functions from
subharmonic functions combines the maximum principle with barriers controlling
boundary behavior. Newtonian-potential estimates for a solution of
$\Delta u=f$ give, for $0<\alpha<1$ and $\Omega'\Subset\Omega$,

\dcref{ch1:1.3}\group{chapter-1}\source{gilbarg-trudinger:page-0015}
\begin{equation}\label{gt.eq1.5}
\tag{1.3}
\|u\|_{C^{2,\alpha}(\bar{\Omega}')} \leq C\!\left(\sup_{\Omega}|u|+\|f\|_{C^{\alpha}(\bar{\Omega})}\right).
\end{equation}

This is the content of Theorems 4.6 and 4.8. Here $C$ depends only on $\alpha$,
$n$, and $\operatorname{dist}(\Omega',\partial\Omega)$. With sufficiently smooth
boundary and boundary data, an analogous estimate holds globally; the Chapter 4
boundary estimates treat hyperplanes and spheres.

Schauder theory extends this estimate to \eqref{gt.eq1.1} with H\"older-continuous
coefficients by freezing the principal coefficients and treating the variable
coefficient operator as a perturbation. The resulting constant also depends on
coefficient bounds and H\"older seminorms and on the extreme eigenvalues of
$[a^{ij}]$. The interior weighted and global estimates are Theorems 6.2 and 6.6;
they yield solvability by the method of continuity (Theorem 6.8) and higher
regularity under stronger smoothness hypotheses (Theorems 6.17 and 6.19).
These are \emph{apriori estimates}: bounds determined by the data and valid for
every possible solution, independently of whether existence has yet been
established. When $c\leq0$, the maximum principle supplies comparison estimates
and uniqueness. The classical existence and regularity conclusions generally
fail when the coefficients are only continuous. The estimates of Chapters 4 and
6 also admit proofs based entirely on maximum-principle comparison, without
Newtonian-potential integrals.

Section 6.5 obtains Dirichlet solvability for continuous boundary data on a broad
class of domains using only interior estimates. Section 6.6 treats certain
non-uniformly elliptic equations, where boundary geometry and degenerating
ellipticity govern attainment of continuous boundary data. Section 6.7 treats the
regular oblique derivative problem.

For coefficients of lower regularity, consider the divergence-form operator

\dcref{ch1:1.6}\group{chapter-1}\source{gilbarg-trudinger:page-0017}
\begin{equation}\label{gt.eq1.6}
L'u \equiv D_i(a^{ij}(x)D_j u + b^i(x)u) + c^i(x)D_i u + d(x)u.
\end{equation}

If $a^{ij}$, $b^i$, and $c^i$ are bounded and measurable and $g$ is integrable,
a function $u\in W^{1,2}(\Omega)$ is a weak solution of $L'u=g$ when

\dcref{ch1:1.4}\group{chapter-1}\source{gilbarg-trudinger:page-0017}
\begin{equation}\label{gt.eq1.7}
\int_{\Omega} \bigl[(a^{ij}D_j u+b^i u)D_i v-(c^i D_i u+d u)v\bigr]\,dx
=-\int_{\Omega} g v\,dx
\end{equation}

for every $v\in C_0^1(\Omega)$. If the coefficients and $g$ are sufficiently
smooth and $u\in C^2(\Omega)$, the weak identity gives the classical equation.
For boundary data $\varphi\in W^{1,2}(\Omega)$, the weak Dirichlet condition is
$u-\varphi\in W_0^{1,2}(\Omega)$. If the least eigenvalue of $[a^{ij}]$ is
bounded away from zero,

\dcref{ch1:1.5}
\group{chapter-1}
\source{gilbarg-trudinger:page-0017}
\begin{equation}\label{gt.eq1.8}
D_i b^i+d\leq 0
\end{equation}

holds weakly, and $g\in L^2(\Omega)$, the generalized Dirichlet problem has a
unique solution in $W^{1,2}(\Omega)$ (Theorem 8.3). The sign condition gives the
weak maximum principle (Theorem 8.1) and uniqueness; existence follows from the
Fredholm alternative (Theorem 8.6) and the Riesz representation theorem on
$W_0^{1,2}(\Omega)$.

Additional coefficient regularity places weak solutions in higher $W^{k,2}$
spaces (Theorems 8.8 and 8.10); Sobolev embedding then makes them classical when
the hypotheses are strong enough. Boundary regularity promotes interior estimates
to global ones. For quasilinear equations, one bootstraps
by substituting the current weak solution into the coefficients of an associated
linear equation, improving its regularity, and iterating. Harnack inequalities
obtained by Moser iteration from integral test-function estimates give the
H\"older bounds used in this procedure (Theorems 8.17, 8.18, 8.20, and 8.24).
The same theory also treats the Wiener criterion for regular boundary points,
eigenvalues and eigenfunctions, and first-derivative H\"older estimates for linear
divergence-form equations.

A strong solution has weak second derivatives and satisfies
\eqref{gt.eq1.1} almost everywhere. The corresponding theory includes the
Aleksandrov maximum principle and estimate in $W^{2,n}(\Omega)$ (Theorem 9.1),
the Krylov--Safonov Harnack and H\"older estimates (Theorems 9.20 and 9.22;
Corollaries 9.24 and 9.25), and the Calder\'on--Zygmund estimate for Poisson's
equation, obtained through Marcinkiewicz interpolation (Theorem 9.9). Interior
and global $W^{2,p}$ estimates for
$1<p<\infty$ (Theorems 9.11 and 9.13) yield the strong Dirichlet theory
(Theorem 9.15 and Corollary 9.18).

\section*{Quasilinear and fully nonlinear theory}

Besides \eqref{gt.eq1.2}, the nonlinear theory treats divergence-form operators

\dcref{ch1:1.6}\group{chapter-1}\source{gilbarg-trudinger:page-0018}
\begin{equation}\label{gt.eq1.9}
\tag{1.6}
Qu \equiv \diver A(x, u, Du) + B(x, u, Du),
\end{equation}

where $A$ is vector-valued and $B$ scalar-valued on
$\Omega\times\R\times\R^n$. Maximum and comparison principles yield apriori
bounds under structure conditions that may be weaker than ellipticity
(Theorem 10.9).

Under the hypotheses of the classical Dirichlet theory, a smooth solution of
$Qu=0$ can be represented as a fixed point $u=Tu$ of a compact map on
$C^{1,\alpha}(\overline\Omega)$, where $Tu$ solves the linear equation obtained by
inserting $u$ into the coefficients. Leray--Schauder reduces existence to a uniform
$C^{1,\alpha}$ bound for solutions of
$u=T(u;\sigma)$, $0\leq\sigma\leq1$, with $T(u;1)=Tu$ (Theorems 11.4 and
11.6). The estimates are obtained successively for
\[
\sup_{\Omega}|u|,\qquad
\sup_{\partial\Omega}|Du|,\qquad
\sup_{\Omega}|Du|,\qquad
\|u\|_{C^{1,\alpha}(\overline\Omega)}.
\]
The first bound is supplied by Chapter 10, assumed, or implied by the equation.

In two dimensions, uniformly elliptic linear equations admit apriori
$C^{1,\alpha}$ estimates depending only on ellipticity constants and coefficient
bounds, without coefficient regularity assumptions; quasiconformal and
divergence-form methods provide the corresponding estimates. For the elliptic
equation

\dcref{ch1:1.7}\group{chapter-1}\source{gilbarg-trudinger:page-0019}
\begin{equation}\label{gt.eq1.10}
au_{xx}+2bu_{xy}+cu_{yy}=0,
\end{equation}

on a bounded convex domain, continuous boundary data satisfying a bounded-slope
or three-point condition with constant $K$ give the apriori estimate $|Du|\leq K$
(Lemma 12.6). This also applies when $a,b,c$ depend on
$(x,y,u,u_x,u_y)$ and reduces the quasilinear Dirichlet problem to the uniformly
elliptic case (Theorem 12.5); local H\"older continuity of the coefficients suffices
for existence with bounded-slope boundary data (Theorem 12.7). The
$C^{1,\alpha}$ estimate above is Theorem 12.4, and the two-dimensional
quasilinear theory otherwise requires only the Schauder fixed point theorem
(Theorem 11.1) from Chapters 7--11.

The higher-dimensional existence scheme uses H\"older estimates for derivatives,
boundary gradient estimates, and global and interior gradient estimates. The
boundary theory treats general and convex domains and relates generalized
curvature conditions to Dirichlet solvability. For the minimal-surface equation
on a smooth domain, solvability for every smooth boundary value is equivalent to
nonnegative boundary mean curvature with respect to the inner normal
(Theorem 14.14). Bernstein-type arguments control the interior gradient by the
boundary gradient for uniformly elliptic equations with natural growth and for
equations with prescribed-mean-curvature structure (Theorem 15.2); a variant
gives interior estimates for a narrower class (Theorem 15.3). Divergence-form
equations admit further gradient estimates under alternative coefficient
conditions (Theorems 15.6, 15.7, and 15.8). Combining these estimates with the
Chapter 10 bounds and the continuation framework of Theorem 11.8 gives the
Chapter 15 existence results.

For the prescribed-mean-curvature equation, hypersurface identities for the
tangential gradient and Laplacian yield interior gradient bounds (Theorem 16.5)
and Dirichlet existence with continuous boundary values (Theorems 16.8 and
16.10). In two variables, generalized quasiconformal methods give interior first-
and second-derivative estimates, including curvature estimates for minimal graphs
(Theorem 16.20) and the conclusion that entire two-dimensional solutions of the
minimal-surface equation are linear (Corollary 16.19).

Fully nonlinear equations have the form

\dcref{ch1:1.8}\group{chapter-1}\source{gilbarg-trudinger:page-0020}
\begin{equation}\label{gt.eq1.11}
F[u] = F(x,u,Du,D^2u) = 0.
\end{equation}

Here $F$ is defined for
$(x,z,p,r)\in\Omega\times\R\times\R^n\times\R^{n\times n}$, with
$\R^{n\times n}$ the symmetric real matrices, and the equation is elliptic when
$F_r$ is positive definite. The method of continuity reduces the Dirichlet
problem to apriori $C^{2,\alpha}(\overline\Omega)$ estimates (Theorem 17.8).
Such estimates are developed for two-dimensional equations (Theorems 17.9 and
17.10), uniformly elliptic equations (Theorems 17.14 and 17.15), and equations
of Monge--Amp\`ere type (Theorems 17.19, 17.20, and 17.26). They yield the
uniformly elliptic existence theorems 17.17 and 17.18 and the Monge--Amp\`ere
existence Theorem 17.23; Bellman--Pucci operators are included in the fully
nonlinear class.

\bigskip

\begin{notation}\label{gt.not.basic-notation}\dcref{ch1:basic-notation}\group{chapter-1}
\source{gilbarg-trudinger:page-0022,gilbarg-trudinger:page-0023}
\noindent\textbf{Basic Notation}

\begin{description}
\item[$\R^n$:] Euclidean $n$-space, $n\geq2$, with
$x=(x_1,\dots,x_n)$ and
$|x|=(\sum_i x_i^2)^{1/2}$. For an ordered $n$-tuple
$b=(b_1,\dots,b_n)$, $\|b\|=(\sum_i b_i^2)^{1/2}$.
\item[$\R^n_+$:] The half-space $\{x\in\R^n:x_n>0\}$.
\item[$\partial S$, $\overline S$:] The boundary of $S$ and its closure
$\overline S=S\cup\partial S$.
\item[$S-S'$:] The set $\{x\in S:x\notin S'\}$.
\item[$S'\Subset S$:] The closure of $S'$ is a compact subset of $S$.
\item[$\Omega$:] A proper open subset of $\R^n$, not necessarily bounded; it is
a domain when connected, and $|\Omega|$ denotes its volume.
\item[$B(y)$, $B_r(y)$:] A ball centered at $y$, and the open ball of radius
$r$ centered at $y$.
\item[$\omega_n$:] The volume of the unit ball in $\R^n$,
$\omega_n=2\pi^{n/2}/(n\Gamma(n/2))$.
\item[$D_i u$, $D_{ij}u$:] The derivatives
$\partial u/\partial x_i$ and
$\partial^2u/(\partial x_i\partial x_j)$; $Du=(D_1u,\dots,D_nu)$ and
$D^2u=[D_{ij}u]$ are the gradient and Hessian.
\item[$\beta$:] A multi-index
$\beta=(\beta_1,\dots,\beta_n)$ with $\beta_i\in\mathbb Z_{\geq0}$ and
$|\beta|=\sum_i\beta_i$ and
\[
D^\beta u=
\frac{\partial^{|\beta|}u}
{\partial x_1^{\beta_1}\cdots\partial x_n^{\beta_n}}.
\]
\item[$C^0(\Omega)$, $C^0(\overline\Omega)$:] Continuous functions on
$\Omega$ and $\overline\Omega$, respectively.
\item[$C^k(\Omega)$:] Functions whose derivatives through order $k$ are
continuous in $\Omega$, for $k\in\mathbb Z_{\geq0}$ or $k=\infty$.
\item[$C^k(\overline\Omega)$:] Functions in $C^k(\Omega)$ whose derivatives
through order $k$ extend continuously to $\overline\Omega$.
\item[$\supp u$:] The closure of $\{x:u(x)\neq0\}$.
\item[$C_0^k(\Omega)$:] Functions in $C^k(\Omega)$ with compact support in
$\Omega$.
\item[$C=C(\cdot,\dots,\cdot)$:] A constant depending only on the displayed
arguments; its value may change between occurrences in the same context.
\end{description}
\end{notation}
